\chapter{Interprocedural and Heap Analysis}
\label{ch:interprocedural-heap-analysis}

A helper that writes through its first argument can be harmless to one caller's cached load and destructive to another's.
The difference is not in the helper's syntax; it is in the objects supplied by the callers and the relationships among their arguments.
Analyzing the helper once is attractive, but a summary must retain enough structure to be instantiated correctly at each call.
Recursion, repeated allocation, and separately compiled code make that structure part of the soundness argument.

Chapter~\ref{ch:ch22-interprocedural-analysis} develops call graphs and may-modify propagation, while Chapter~\ref{ch:ch20-alias-analysis} develops inclusion-based points-to sets.
Here we connect those mechanisms to matched call/return semantics, relational summaries, heap updates, and module interfaces.
The order remains inclusion of represented behavior: $\bot$ means no behavior, $\top$ means unconstrained behavior, and joins cover all alternatives.

\section{A call has a matching continuation}
\label{sec:ch44-matched}

Consider two calls to a pure helper that returns its argument plus two.
Call site $c_1$ supplies $4$ and receives $6$; site $c_2$ supplies $9$ and receives $11$.
If an interprocedural graph connects both call sites to one entry and its exit to both continuations, ordinary graph reachability admits an impossible path: enter from $c_1$ and return to the continuation of $c_2$.
Concrete execution maintains a return stack, so it cannot take that path.

Figure~\ref{fig:ch44-matched-returns} separates shared code from the two legal return connections.
A \term{matched call/return path} returns from each invocation to the continuation belonging to that invocation.
Nested calls must obey the same nesting discipline.
Exceptions, tail calls, continuations, and coroutine transfers need their own control rules; the baseline here is ordinary sequential call and return.
Ignoring those distinctions can invent behaviors and weaken analysis, while omitting legal exceptional returns can be unsound.

\begin{figure}
\centering
\begin{tikzpicture}[node distance=12mm and 15mm]
\node[basic block] (c1) {$c_1$: argument 4};
\node[basic block,right=of c1] (c2) {$c_2$: argument 9};
\node[basic block,below right=15mm and -4mm of c1] (helper) {shared helper\\return argument + 2};
\node[basic block,below=40mm of c1] (r1) {$r_1$: result 6};
\node[basic block,below=40mm of c2] (r2) {$r_2$: result 11};
\draw[cfg edge] (c1) -- node[left,font=\small] {$c_1$} (helper);
\draw[cfg edge] (c2) -- node[right,font=\small] {$c_2$} (helper);
\draw[cfg edge] (helper) -- node[left,font=\small] {match $c_1$} (r1);
\draw[cfg edge] (helper) -- node[right,font=\small] {match $c_2$} (r2);
\end{tikzpicture}
\caption{The body is shared, but each summary application returns to its own continuation. A path entering at $c_1$ cannot select $r_2$.}
\label{fig:ch44-matched-returns}
\end{figure}


A procedure summary avoids traversing the callee's graph separately for every dynamic call.
For a pure helper, the relation $r=a+2$ links entry argument $a$ to normal result $r$.
Instantiating it at $c_1$ yields $r=6$; instantiating the same relation at $c_2$ yields $r=11$.
This is more precise than storing only the union of observed return values $\{6,11\}$.
A shared summary need not be context-insensitive in the sense of erasing all input/output correlation.

At a call, save the caller-local facts that the callee cannot modify, bind actual arguments to fresh formal names, apply the summary, and map the result back to the designated continuation.
All fresh names must be renamed apart from names belonging to another invocation.
The caller's untouched state is the frame around the call.
Memory in that frame can be retained only when the summary's effect contract establishes that it is untouched.

\section{Specify summaries as relations and effects}
\label{sec:ch44-summary-contract}

For a sequential procedure $P$, its normal-return semantics can be described by a relation between entry and exit states.
An entry state contains arguments and the relevant pre-call heap; an exit state contains the result and the post-call heap.
A sound relational summary covers every such terminating input/output pair.
The summary need not assert that every allowed input terminates.
In particular, a relation containing no output for an input is not, by itself, a proof that executing the call has no effects.

For general compiler use, separate normal returns, exceptional returns, and effects possible before either return or divergence.
A procedure that stores to a global and then loops forever has no normal-return pair, but the store remains observable.
A normal-return relation can justify facts after a call; it cannot alone justify removing the call.
Likewise, a may-write set does not prove that a location is overwritten on every return path, and a no-write summary does not prove absence of reads.

Parameterize heap effects by formal roots when possible.
A summary saying that the procedure may write field $f$ of the object reached by argument one can be instantiated with the caller's points-to information.
Replacing that symbolic root by the union of every object ever passed to the helper loses precision at every call.
The instantiation must also account for aliases: formal arguments one and two may denote the same object.
The same heap cannot be modeled as two independent copies merely because the interface has two parameter names.

A compact summary might contain a precondition, a return relation, may-read and may-write footprints, escape information, and exceptional outcomes.
Each component has a separate interpretation.
The precondition states when the more precise relation is valid; it is not a fact the caller may assume without proof.
If the caller cannot establish it, use a weaker unconditional summary or retain a runtime guard.
The footprint states which cells may be touched, not which cells definitely are touched.
Escape information states which references may outlive the call, affecting later allocation and lifetime reasoning.

\section{Contexts preserve selected distinctions}
\label{sec:ch44-contexts}

A \term{call string} records a sequence of call-site identifiers used to distinguish analysis contexts.
Keeping only the most recent $k$ sites gives a finite policy for a finite program.
Different invocations with the same retained suffix are merged, including sufficiently deep recursive invocations.
The empty suffix is the context-insensitive case.
Context keys must be bounded deliberately; the dynamic call stack is unbounded even when the source has few procedures.

An alternative key is an abstract input description, such as a constant argument, a range class, or a selected alias partition.
Memoizing by exact interval endpoints without a merge policy can create unboundedly many contexts.
When a budget is reached, merge inputs and recompute a summary covering the merged description, or use a previously established general summary.
Dropping the expensive context would drop behaviors.

Context sensitivity does not replace matched returns.
A bounded call string identifies an abstract invocation class, while the analysis still needs to route a callee result to a call site's own continuation.
Even when several callers share one context, each summary application is attached to its calling edge.
Conversely, correctly matched summaries can still lose precision when their input/output abstraction is too coarse.
These are two different design decisions~\cite{nielson-1999-principles}.

For a receiver-oriented language, an allocation-site or object-based context can distinguish methods operating on different receiver classes.
Such a context distinguishes abstract objects, not necessarily dynamic instances.
One allocation site in a loop can still represent many simultaneous receivers.
The heap model must decide what identity information survives inside the context.

\section{Recursive summaries are a fixed-point problem}
\label{sec:ch44-recursion}

A recursive procedure's summary depends on itself, directly or through other procedures.
For a fully specified baseline, assume a finite abstract boundary-state universe $Q$ with $m=|Q|$ states.
Local abstract transitions form finite relations and soundly simulate concrete steps.
Call binding, return projection, and preservation of caller-local state are included in the model.
If the concrete language is infinite-state, constructing this finite model is itself an abstraction obligation.

Store a summary $S_P\subseteq Q\times Q$ for each procedure $P$.
It covers normal terminating entry/exit pairs.
The lattice order is subset inclusion and join is union.
Let $\operatorname{Eval}(P,S)$ compute intraprocedural reachability while replacing a call to $R$ by the currently stored relation $S_R$, applied at that call's matching continuation.
Retain the entry state as an origin tag throughout this local computation, so that outputs are paired with their own inputs.
Loops inside a body are solved by finite reachability on these tagged states.

Algorithm~\ref{alg:ch44-summary-fixedpoint} iterates all summaries synchronously.
The local evaluator is monotone in the summaries it reads.
Bottom initialization initially describes only behavior with no completed calls; later rounds add behavior justified by increasingly deep finite call trees.
Every terminating execution has a finite call tree, so the least fixed point includes all terminating modeled behavior.
Intermediate summaries are incomplete and must not be exposed as sound whole-program answers before convergence.

\begin{algorithm}
\caption{Compute finite relational summaries for mutually recursive procedures.}
\label{alg:ch44-summary-fixedpoint}
\begin{pseudocode}
\Require finite procedures, boundary universe $Q$, sound finite local model
\Require $\operatorname{Eval}$ computes tagged local reachability with matched calls
\Ensure $S_P$ covers every normal terminating modeled pair for each $P$
\State $S_P\gets\emptyset$ for every procedure $P$
\Loop
  \ForAll{procedures $P$}
    \State $T_P\gets S_P\cup\operatorname{Eval}(P,S)$
  \EndFor
  \If{$T_P=S_P$ for every $P$}
    \State \Return $S$
  \EndIf
  \State $S\gets T$
\EndLoop
\end{pseudocode}
\end{algorithm}

Let $p$ count procedures, $N$ all intraprocedural points, $L$ ordinary edges, and $C$ call sites with their possible target edges.
There are at most $pm^2$ summary memberships to add, hence at most $pm^2+1$ complete sweeps including the final unchanged sweep.
One straightforward evaluator explores up to $m^2$ origin/current-state pairs per program point.
For each origin, a dense local relation or callee relation can have $m^2$ transition pairs.
Thus a whole sweep can be implemented in $O(Nm^2+(L+C)m^3)$ membership and transition operations, assuming constant-time table access.
The resulting coarse bound is the sweep cost multiplied by $pm^2+1$; storage is $O((N+p)m^2)$ cells plus the model's transition relations.
These are bounds for the explicitly tabulated finite model, not for symbolic summaries whose application may invoke a solver.

Production solvers avoid these full sweeps.
They process strongly connected call-graph components, register summary readers, and propagate changed pairs or changed abstract values.
Acyclic components can be handled after their callees, while recursive components require iteration.
For infinite-height numerical summaries, widening must also occur at recursive summary dependencies, not only at intraprocedural loop headers.
A program can recurse indefinitely without containing a CFG back edge inside any one procedure.

\section{Compute a recursive input/output relation}
\label{sec:ch44-recursive-example}

Consider a helper with precondition $0\leq n\leq3$.
It returns $2$ when $n=0$; otherwise it returns one plus the result of calling itself on $n-1$.
Use mathematical integers and normal returns only.
The recursive call remains inside the precondition, and the result stays between two and five.
The finite example can therefore be represented exactly with boundary values from zero through five, restricting admissible entries to zero through three.

Let $S^{(0)}=\emptyset$.
In the first evaluation, only the base case yields a return, so $S^{(1)}=\{(0,2)\}$.
The second evaluation can use that pair at the recursive call with argument zero and obtains $(1,3)$.
Successive rounds add $(2,4)$ and $(3,5)$.
Table~\ref{tab:ch44-recursive-pairs} records every change, including the unchanged convergence check.

\begin{table}
\centering
\caption{Each synchronous round adds one deeper terminating invocation to the summary.}
\label{tab:ch44-recursive-pairs}
\begin{tabular}{p{0.16\linewidth}p{0.28\linewidth}p{0.45\linewidth}}
\toprule
Round & Newly established pair & Entire normal-return relation \\
\midrule
0 & None & $\emptyset$ \\
1 & $(0,2)$ & $\{(0,2)\}$ \\
2 & $(1,3)$ & $\{(0,2),(1,3)\}$ \\
3 & $(2,4)$ & $\{(0,2),(1,3),(2,4)\}$ \\
4 & $(3,5)$ & $\{(0,2),(1,3),(2,4),(3,5)\}$ \\
5 & None & Same as round 4 \\
\bottomrule
\end{tabular}
\end{table}

At convergence, a call with argument one returns three and a call with argument three returns five.
Projecting the relation to an output-only set $\{2,3,4,5\}$ would lose that distinction.
An interval summary retaining only input $[0,3]$ and output $[2,5]$ loses it too.
A relational numerical summary $r=n+2$ could express the same input/output dependence compactly and cover a larger range, provided its arithmetic and recursive-invariant proof were established.

This example also shows why a recursive call cannot be assigned an empty summary as a final conservative fallback.
Empty means no normal behavior, not unknown normal behavior.
During least-fixed-point construction it is a seed; at a timeout it remains incomplete.
A compiler that cannot finish must use a separately sound cover, such as an unconstrained return with conservative effects, rather than return the partial relation to clients.

\section{Heap names are abstractions of objects}
\label{sec:ch44-heap-names}

An \term{allocation-site abstraction} names heap objects by the instruction that allocates them, optionally combined with a bounded context.
A field-sensitive cell is then a pair $(a,f)$ of an abstract object $a$ and a field $f$.
The heap map $H^\#$ assigns an abstract value to each such cell.
Its concretization requires every concrete cell represented by $(a,f)$ to have a value allowed by $H^\#(a,f)$.
If several concrete objects share the name $a$, the abstract value must cover all of their fields.

Figure~\ref{fig:ch44-summary-heap} shows two live allocations from one site.
Their concrete fields initially contain zero.
Pointers $p$ and $q$ refer to different objects, but the points-to abstraction gives both the singleton set $\{a\}$.
That equality does not prove $p=q$.
The singleton set contains one abstract name, and the name represents two concrete objects.

\begin{figure}
\centering
\begin{tikzpicture}[node distance=13mm and 17mm]
\node[basic block] (p) {$p$};
\node[basic block,below=of p] (q) {$q$};
\node[basic block,right=of p] (o1) {object 1: f=0};
\node[basic block,right=of q] (o2) {object 2: f=0};
\node[basic block,right=20mm of o1,yshift=-12mm] (site) {site a\\f in \{0\}};
\draw[cfg edge] (p) -- (o1);
\draw[cfg edge] (q) -- (o2);
\draw[cfg edge,dashed] (o1) -- (site);
\draw[cfg edge,dashed] (o2) -- (site);
\node[font=\small,below=7mm of o2] {two live objects, one abstract name};
\end{tikzpicture}
\caption{Dashed abstraction arrows merge two objects allocated at the same site. A store through $p$ must retain the old value of the object reached by $q$.}
\label{fig:ch44-summary-heap}
\end{figure}


An abstract name can instead be known to represent at most one live object per concrete state at a particular program point.
For example, a distinguished fresh allocation can remain separate from older objects until another allocation at the same site forces a merge.
The uniqueness fact is temporal and state-dependent.
Allocation in a loop, recursive activations, escaping references, and joins can invalidate it.
It is not enough to inspect the allocation instruction and notice that it appears only once in the source.

Heap abstractions also need a policy for external objects, null, invalid addresses, and deallocated objects.
An empty points-to set must not stand for an arbitrary foreign pointer.
A sound unknown-address element or an explicit universe of possible external regions gives unknown code somewhere to point.
Lifetime and invalid-access outcomes must remain represented even when the compiler is interested only in ordinary load values.

\section{Strong updates require a singleton concrete target}
\label{sec:ch44-updates}

A \term{strong update} replaces a cell's old abstract value with the value being stored.
A \term{weak update} joins the stored value with the old value.
The distinction follows from the heap concretization.
If an abstract cell represents several concrete cells and the store changes only one, the others retain their old contents and must remain covered.

Algorithm~\ref{alg:ch44-store} specifies a field store for a restricted sequential heap model.
Fields are distinct nonoverlapping cells, addresses have already been checked valid on the normal path, and a store writes one entire field with compatible representation.
The target set is a sound nonempty set of possible abstract field cells.
The predicate $\operatorname{One}(a,f)$ proves that this cell represents exactly one concrete target in every represented state and that the store necessarily writes it.
For a singleton target set, this includes the requisite object multiplicity and offset facts.

\begin{algorithm}
\caption{Update an abstract heap for a valid whole-field store.}
\label{alg:ch44-store}
\begin{pseudocode}
\Require heap $H$, nonempty sound target set $T$, stored abstract value $v$
\Require nonoverlapping fields; sound singleton-target predicate $\operatorname{One}$
\Ensure returned heap covers every normal concrete store result
\State $H'\gets$ a copy of $H$
\If{$T=\{(a,f)\}$ and $\operatorname{One}(a,f)$}
  \State $H'(a,f)\gets v$
\Else
  \ForAll{$(a,f)\in T$}
    \State $H'(a,f)\gets H(a,f)\join v$
  \EndFor
\EndIf
\State \Return $H'$
\end{pseudocode}
\end{algorithm}

Other pointers may alias the target of a strong update.
The algorithm does not require that only one pointer names the cell; it requires that the store definitely changes the one represented concrete cell.
Because all aliases read the shared heap map, they observe the replaced value.
An implementation that caches load values separately must invalidate those aliases' cached facts as well.

For the two allocations in Figure~\ref{fig:ch44-summary-heap}, storing seven through $p$ changes the first concrete field from zero to seven and leaves the second at zero.
The site cell must therefore become $\{0,7\}$ in a finite value-set domain, or $[0,7]$ in an interval domain.
Replacing it by $\{7\}$ would misdescribe the subsequent load through $q$.
Table~\ref{tab:ch44-updates} computes this case and two contrasting cases.

\begin{table}
\centering
\caption{The update strength depends on concrete multiplicity and definite targeting, not only the number of abstract names.}
\label{tab:ch44-updates}
\begin{tabular}{p{0.28\linewidth}p{0.27\linewidth}p{0.33\linewidth}}
\toprule
Target information & Old field values & Result of store \\
\midrule
One name $a$, two live objects & $H(a,f)=\{0\}$ & Store $7$: $H'(a,f)=\{0,7\}$ \\
One unique target $u$ & $H(u,f)=\{2\}$ & Store $5$: $H'(u,f)=\{5\}$ \\
Either unique $u$ or unique $v$ & $H(u,f)=\{2\}$, $H(v,f)=\{9\}$ & Store $5$: $H'(u,f)=\{2,5\}$, $H'(v,f)=\{5,9\}$ \\
\bottomrule
\end{tabular}
\end{table}

In the final row, each object is individually unique, but the analysis does not know which one the store targets.
Each cell may remain unchanged, so both updates are weak.
A path partition that proves the target on each arm can strongly update within each partition and then join.
The joined per-cell sets are the same here, but the partitions retain correlations between the pointer and the changed cell that a product heap loses.

If $A$ abstract cells are stored in a flat map and joining values costs $J$, the displayed copying implementation costs $O(A+|T|J)$ time and $O(A)$ new map storage per store.
A persistent or copy-on-write map changes that cost; the join itself is not always constant time.
For value sets over $U$ possible values stored as bitsets, $J=O(\lceil U/w\rceil)$ word operations at word width $w$.
The cost of proving $\operatorname{One}$ and computing $T$ belongs to the pointer and multiplicity analyses and must be added.

\section{Aliasing connects summary application to memory optimization}
\label{sec:ch44-alias-summary}

Suppose a helper writes five through parameter $p$ and returns the field read through parameter $q$.
If the caller proves that $p$ and $q$ select disjoint cells, the returned value can retain the old fact for $q$.
If they must select the same unique cell, the return is five.
If they may alias, the return must cover both cases.
A summary that assumes the formal parameters are distinct without an explicit precondition is unsound.

One design keeps separate alias cases in the summary and chooses or joins cases at instantiation.
Another uses symbolic heap roots and lets substitution identify equal actual roots.
Either way, aliasing must be resolved before relying on frame preservation.
Taking a summary's may-write footprint and subtracting it syntactically from the caller's known locations is insufficient when two different names can overlap.

Real memory accesses name byte ranges, not only fields.
Partial stores, unions, packed representations, and type-punning can invalidate the nonoverlapping-field premise of Algorithm~\ref{alg:ch44-store}.
A byte-aware implementation updates every overlapping abstract cell or uses a coarser region abstraction.
Equal starting addresses do not establish compatible load/store sizes, and disjoint field names do not establish disjoint bytes in every language layout.
The abstract heap model must match the IR memory model consumed by the optimizer.

Dynamic call targets add another dependency.
A points-to result can discover a new possible callee, which enlarges the summary applied at the call and can in turn enlarge points-to sets.
An integrated solver must propagate both kinds of change to a joint fixed point, or use a staged conservative call graph whose target set already covers every possible call.
Using an incomplete call graph while waiting for more precise pointer information is unsafe if clients observe the intermediate summaries.

\section{Module interfaces carry assumptions across builds}
\label{sec:ch44-modules}

Separate compilation replaces unseen bodies with interface contracts.
A useful serialized summary identifies the function, calling convention, semantic version of the summary format, target data layout, and arithmetic and memory assumptions.
It also records the body or interface revision on which it depends.
Two summaries with the same source-level name need not describe the same implementation after linking, symbol interposition, or dynamic loading.
The optimizer must know which binding assumptions the build permits.

An unavailable body does not justify an empty effect set.
Without a stronger interface, an external call may access reachable mutable memory and relevant globals, retain references, invoke permitted callbacks, and return any value allowed by its type and language contract.
A local object proved inaccessible to the callee can remain outside this footprint.
That exclusion depends on escape and reachability facts, not merely on the object's storage class.

Callbacks make effects transitive in both directions across an interface.
A library routine that performs no direct store may still invoke a client callback that mutates the caller's global state.
The interface must expose the callback possibility or conservatively include its effects.
Similarly, a no-capture promise affects later calls: if it is absent, a pointer passed now may become reachable by unknown code later.
The call's immediate may-write set is not enough to reason about that future access.

Summary reuse requires invalidation when any assumption changes.
Changing a callee body invalidates its summary and readers; changing the alias policy, data layout, or external binding assumption can invalidate an entire family of summaries.
The dependency traversal from Chapter~\ref{ch:designing-sound-analyses} applies here, but the dependency keys include semantic options and imported interfaces.
A version mismatch should trigger recomputation or a conservative fallback, not reinterpret old bytes under a new meaning.

For compiler clients, expose only guarantees whose quantifiers are clear.
\enquote{May write $a$} includes all possible writes but promises none.
\enquote{Returns a fresh object} requires identity and lifetime facts, not merely an allocation-site label.
\enquote{Argument is not captured} must specify whether callbacks or exceptional returns are included.
Well-specified interfaces let separate compilation retain useful facts without treating unavailable code as if it were absent.

\section{Exercises}
\label{sec:ch44-exercises}

\begin{enumerate}
\item \label{ex:ch44-matching}
A helper returns its argument plus two and is called with $1$ at site $a$ and $8$ at site $b$.
Give the two correct return values and one impossible call/return combination admitted by an unmatched graph.
Explain how an input/output relation avoids merging the answers.
\item \label{ex:ch44-recursion}
Change the recursive helper's precondition to $0\leq n\leq4$.
Starting from the empty summary, list the new pair in each changing synchronous round and the number of the first unchanged round.
\item \label{ex:ch44-weak}
Three live objects allocated at one site have field values $1,1,1$.
A store changes exactly one of them to $6$.
Compute the site cell in a finite value-set domain and in an interval domain.
Why does a singleton points-to set not justify replacing the site cell by $6$?
\item \label{ex:ch44-alias-case}
Unique cells $u$ and $v$ initially contain $2$ and $9$.
A helper stores $5$ through $p$ and returns the load through $q$.
Compute its result when $(p,q)=(u,v)$, when $(p,q)=(u,u)$, and when $p=u$ while $q$ may denote either cell.
State the alias precondition needed to reuse the old value of $q$.
\end{enumerate}

\section{Further reading}
\label{sec:ch44-reading}

Nielson, Nielson, and Hankin develop interprocedural paths and context-sensitive analyses in a framework that makes the underlying semantic distinctions explicit~\cite{nielson-1999-principles}.
Cousot and Cousot provide the fixed-point abstraction principles needed when recursive summaries or heap descriptions use infinite abstract domains~\cite{cousot-1977-abstract}.
Kildall's dataflow framework remains useful for understanding the propagation engine underneath finite summary computations~\cite{kildall-1973-dataflow}.
